% !TEX root = main.tex
\chapter{Reasonable Ordered Expansions}
\label{chap:reasonable-expansions}

In the previous chapters, ordered Fraïssé classes were used to obtain extremely amenable automorphism groups and to prepare the description of universal minimal flows. To pass from these ordered expansions back to the underlying unordered structures, one must understand how an ordered class is related to the class obtained by forgetting the linear order. The aim of this chapter is to isolate the conditions under which this passage behaves well.

We first introduce the notion of a reasonable ordered expansion and show that, for a Fraïssé order class \(\cK\), reasonableness is exactly the condition ensuring that the reduct class \(\cK_0\) is again a Fraïssé class and that the reduct of the Fraïssé limit of \(\cK\) is the Fraïssé limit of \(\cK_0\). We then discuss freely ordered expansions and strong amalgamation, and conclude with a brief remark on order-forgetful expansions, where Ramsey properties pass directly between the ordered and unordered settings.

\section{Reducts and expansions}

Let \(L\) be a signature containing a distinguished binary relation symbol \(<\), and let \(L_0=L\setminus\{<\}\). If \(\bfa\) is an \(L\)-structure with underlying set \(A\), we write \(\bfa_0=\bfa\!\restriction\!L_0\) for the reduct of \(\bfa\) to the signature \(L_0\). Thus \(\bfa_0\) is obtained from \(\bfa\) by forgetting the interpretation of the symbol \(<\). Conversely, if \(\bfa_0\) is an \(L_0\)-structure and \(\prec^{\bfa}\) is a linear ordering of \(A\), we write \(\langle \bfa_0,\prec^{\bfa}\rangle\) for the corresponding \(L\)-expansion of \(\bfa_0\) in which \(<\) is interpreted by \(\prec^{\bfa}\).

If \(\cK\) is a class of finite \(L\)-structures, we write \(\cK_0=\{\bfa_0 : \bfa\in\cK\}\) for the class of all reducts of members of \(\cK\) to the signature \(L_0\). In the situations that matter for us, \(\cK\) will be an order class, so each structure in \(\cK\) carries a linear ordering, while \(\cK_0\) is the corresponding unordered class obtained by forgetting that ordering.

The basic question is then the following. Suppose that \(\cK\) is a Fraïssé order class with Fraïssé limit \(\bff\). Under what conditions is the reduct class \(\cK_0\) again a Fraïssé class, and when is its Fraïssé limit simply the reduct \(\bff_0=\bff\!\restriction\!L_0\)?
The key condition is the notion of a reasonable ordered expansion, introduced in the next section.

\section{Reasonable Fraïssé order classes}

We now isolate the condition that controls when a Fraïssé order class behaves well after the linear order is forgotten. Intuitively, the issue is whether embeddings in the reduct can be lifted to embeddings in the ordered class by choosing a suitable ordering on the larger structure. This leads to the following definition.

\begin{definition}
    \label{def:reasonable}
    Let \(L\) be a signature with \(L\supseteq\{<\}\), and put \(L_0=L\setminus\{<\}\). Let \(\cK\) be a class of finite \(L\)-structures, and let \(\cK_0=\{\bfa_0 : \bfa\in\cK\}\). We say that \(\cK\) is \emph{reasonable} if whenever \(\bfa_0,\bfb_0\in\cK_0\), \(\pi\colon \bfa_0\hookrightarrow \bfb_0\) is an embedding, and \(\prec^{\bfa}\) is a linear ordering of \(A\) such that \(\bfa=\langle \bfa_0,\prec^{\bfa}\rangle\in\cK\), there exists a linear ordering \(\prec^{\bfb}\) of \(B\) such that \(\bfb=\langle \bfb_0,\prec^{\bfb}\rangle\in\cK\) and \(\pi\colon \bfa\hookrightarrow \bfb\) is an embedding.
\end{definition}

In other words, \(\cK\) is reasonable if every embedding in the reduct can be made compatible with the order after choosing a suitable ordering on the target structure.

The importance of this condition is that it exactly characterises when the reduct of a Fraïssé order class is again a Fraïssé class, with the expected Fraïssé limit.

\begin{proposition}
    \label{prop:reasonable-characterization}
    Let \(L\supseteq\{<\}\) be a signature, let \(L_0=L\setminus\{<\}\), and let \(\cK\) be a Fraïssé order class of finite \(L\)-structures. Put
    \[
        \cK_0=\{\bfa_0 : \bfa\in\cK\},
        \qquad
        \bff=\Flim(\cK),
        \qquad
        \bff_0=\bff\!\restriction\!L_0.
    \]
    Then the following are equivalent.
    \begin{enumerate}
        \item[(i)] \(\cK_0\) is a Fraïssé class and
              \[
                  \bff_0=\Flim(\cK_0).
              \]
        \item[(ii)] \(\cK\) is reasonable.
    \end{enumerate}
\end{proposition}

\begin{proof}
    Assume first that \(\cK_0\) is a Fraïssé class and that \(\bff_0=\Flim(\cK_0)\).
    We show that \(\cK\) is reasonable.

    Let \(\bfa_0,\bfb_0\in\cK_0\), let \(\pi\colon \bfa_0\hookrightarrow \bfb_0\) be an embedding, and let \(\prec^{\bfa}\) be a linear ordering of \(A\) such that \(\bfa=\langle \bfa_0,\prec^{\bfa}\rangle\in\cK\). Since \(\bff=\Flim(\cK)\), there is an embedding \(\varphi\colon \bfa\hookrightarrow \bff\). Forgetting the order, we may also regard \(\varphi\) as an embedding \(\varphi\colon \bfa_0\hookrightarrow \bff_0\). Because \(\bff_0\) is the Fraïssé limit of \(\cK_0\), it has the extension property. Hence there is an embedding \(\psi\colon \bfb_0\hookrightarrow \bff_0\) such that \(\psi\circ\pi=\varphi\).
    Now let \(\prec^{\bfb}\) be the linear ordering on \(B\) pulled back from \(\bff\) by \(\psi\), that is,
    \[
        x\prec^{\bfb} y
        \quad\Longleftrightarrow\quad
        \psi(x)\prec^{\bff}\psi(y).
    \]
    If we put \(\bfb=\langle \bfb_0,\prec^{\bfb}\rangle\), then \(\bfb\) is isomorphic to the substructure of \(\bff\) induced on \(\psi(B)\), so \(\bfb\in\Age(\bff)=\cK\). Moreover, since \(\psi\circ\pi=\varphi\) and \(\varphi\) preserves \(\prec^{\bfa}\), the map \(\pi\colon \bfa\hookrightarrow \bfb\) is an embedding. Thus \(\cK\) is reasonable.

    Conversely, assume that \(\cK\) is reasonable. We show first that \(\cK_0\) is a Fraïssé class.

    Since \(\cK_0=\Age(\bff_0)\),
    the class \(\cK_0\) is hereditary, countable up to isomorphism, and has structures of arbitrarily large cardinality. It therefore remains only to verify the amalgamation property.

    Let \(\bfa_0,\bfb_0,\bfc_0\in\cK_0\), and let \(f\colon \bfa_0\hookrightarrow \bfb_0\) and \(g\colon \bfa_0\hookrightarrow \bfc_0\) be embeddings. Choose a linear ordering \(\prec^{\bfa}\) on \(A\) such that \(\bfa=\langle \bfa_0,\prec^{\bfa}\rangle\in\cK\).
    Since \(\cK\) is reasonable, there exist linear orderings \(\prec^{\bfb}\) and \(\prec^{\bfc}\) on \(B\) and \(C\) such that \(\bfb=\langle \bfb_0,\prec^{\bfb}\rangle\in\cK\) and \(\bfc=\langle \bfc_0,\prec^{\bfc}\rangle\in\cK\), and both \(f\colon \bfa\hookrightarrow \bfb\) and \(g\colon \bfa\hookrightarrow \bfc\) are embeddings. Since \(\cK\) is a Fraïssé class, it has the amalgamation property. Hence there exist \(\bfd\in\cK\) and embeddings \(r\colon \bfb\hookrightarrow \bfd\) and \(s\colon \bfc\hookrightarrow \bfd\) such that \(r\circ f=s\circ g\). Passing to reducts, if \(\bfd_0=\bfd\!\restriction\!L_0\), then \(r\) and \(s\) are also embeddings \(r\colon \bfb_0\hookrightarrow \bfd_0\) and \(s\colon \bfc_0\hookrightarrow \bfd_0\), and still satisfy \(r\circ f=s\circ g\).
    Thus \(\cK_0\) has the amalgamation property, and therefore \(\cK_0\) is a Fraïssé class.

    It remains to show that
    \[
        \bff_0=\Flim(\cK_0).
    \]
    Since \(\Age(\bff_0)=\cK_0\),
    it is enough to verify that \(\bff_0\) is ultrahomogeneous, or equivalently that it has the extension property.

    Let \(\bfa_0,\bfb_0\in\cK_0\), let \(\pi\colon \bfa_0\hookrightarrow \bfb_0\) be an embedding, and let \(\varphi\colon \bfa_0\hookrightarrow \bff_0\) be an embedding. Pull back the order on \(\bff\) along \(\varphi\) and define a linear ordering \(\prec^{\bfa}\) on \(A\) by
    \[
        x\prec^{\bfa} y
        \quad\Longleftrightarrow\quad
        \varphi(x)\prec^{\bff}\varphi(y).
    \]
    Then \(\bfa=\langle \bfa_0,\prec^{\bfa}\rangle\in\cK\), and \(\varphi\) is in fact an embedding \(\varphi\colon \bfa\hookrightarrow \bff\).
    Since \(\cK\) is reasonable, there exists a linear ordering \(\prec^{\bfb}\) on \(B\) such that \(\bfb=\langle \bfb_0,\prec^{\bfb}\rangle\in\cK\) and \(\pi\colon \bfa\hookrightarrow \bfb\) is an embedding. Because \(\bff=\Flim(\cK)\), it has the extension property, so there is an embedding \(\psi\colon \bfb\hookrightarrow \bff\) such that \(\psi\circ\pi=\varphi\). Forgetting the order, \(\psi\) is also an embedding \(\psi\colon \bfb_0\hookrightarrow \bff_0\) with \(\psi\circ\pi=\varphi\).
    Thus \(\bff_0\) has the extension property. It follows that \(\bff_0\) is ultrahomogeneous and \(\Age(\bff_0)=\cK_0\),
    so
    \[
        \bff_0=\Flim(\cK_0).
    \]
\end{proof}
The proposition shows that for Fraïssé order classes, reasonableness is exactly the condition that allows one to pass from the ordered class back to its reduct without losing the Fraïssé structure. In the next section we consider an important source of reasonable order classes, namely freely ordered expansions.



\section[Freely ordered expansions and strong amalgamation]{Freely ordered expansions and\\strong amalgamation}

A particularly important source of ordered expansions is obtained by imposing no compatibility condition between the additional linear order and the underlying \(L_0\)-structure.

\begin{definition}
    Let \(\cK_0\) be a class of finite \(L_0\)-structures. We define
    \[
        \cK_0 * \cLO
        =
        \{\langle \bfa_0,\prec^{\bfa}\rangle : \bfa_0\in\cK_0 \text{ and } \prec^{\bfa} \text{ is a linear ordering of } A\}.
    \]
    We call \(\cK_0 * \cLO\) the \emph{freely ordered expansion} of \(\cK_0\).
\end{definition}

The point is that in a freely ordered expansion every linear order is allowed. In particular, such classes are automatically reasonable.

\begin{lemma}
    Let \(\cK=\cK_0 * \cLO\). Then \(\cK\) is reasonable.
\end{lemma}

\begin{proof}
    Let \(\bfa_0,\bfb_0\in\cK_0\), let \(\pi\colon \bfa_0\hookrightarrow \bfb_0\) be an embedding, and suppose that \(\bfa=\langle \bfa_0,\prec^{\bfa}\rangle\in\cK\).
    Choose any linear ordering \(\prec^{\bfb}\) of \(B\) extending the order transported from \(\prec^{\bfa}\) along \(\pi\). Then \(\bfb=\langle \bfb_0,\prec^{\bfb}\rangle\in\cK\),
    and \(\pi\colon \bfa\hookrightarrow \bfb\) is an embedding.
\end{proof}

For freely ordered expansions, amalgamation in the ordered class is controlled exactly by strong amalgamation in the reduct.
\begin{proposition}
    \label{prop:free-orders-sap}
    Let \(L\supseteq\{<\}\) be a signature, let \(L_0=L\setminus\{<\}\), let \(\cK_0\) be a class of finite \(L_0\)-structures, and put \(\cK=\cK_0 * \cLO\).
    Then the following are equivalent.
    \begin{enumerate}
        \item[(i)] \(\cK\) has the amalgamation property.
        \item[(ii)] \(\cK\) has the strong amalgamation property.
        \item[(iii)] \(\cK_0\) has the strong amalgamation property.
    \end{enumerate}
\end{proposition}

\begin{proof}
    The implication (ii) \(\Rightarrow\) (i) is immediate.

    We prove (i) \(\Rightarrow\) (iii). Let \(\bfa_0,\bfb_0,\bfc_0\in\cK_0\), and let \(f\colon \bfa_0\hookrightarrow \bfb_0\) and \(g\colon \bfa_0\hookrightarrow \bfc_0\) be embeddings. Choose a linear ordering \(\prec^{\bfa}\) on \(A\). Since \(\cK=\cK_0 * \cLO\),
    we may choose linear orderings \(\prec^{\bfb}\) on \(B\) and \(\prec^{\bfc}\) on \(C\) such that
    \[
        \bfa=\langle \bfa_0,\prec^{\bfa}\rangle,\qquad
        \bfb=\langle \bfb_0,\prec^{\bfb}\rangle,\qquad
        \bfc=\langle \bfc_0,\prec^{\bfc}\rangle
    \]
    all belong to \(\cK\), the maps \(f\colon \bfa\hookrightarrow \bfb\) and \(g\colon \bfa\hookrightarrow \bfc\) are embeddings, and in addition
    \[
        f(A)\prec^{\bfb} (B\setminus f(A))
        \qquad\text{and}\qquad
        (C\setminus g(A))\prec^{\bfc} g(A).
    \]
    Since \(\cK\) has the amalgamation property, there exist \(\bfd\in\cK\) and embeddings \(r\colon \bfb\hookrightarrow \bfd\) and \(s\colon \bfc\hookrightarrow \bfd\) such that \(r\circ f=s\circ g\).
    We claim that
    \[
        r(B)\cap s(C)=r(f(A)).
    \]
    Suppose, towards a contradiction, that there is \(d\in r(B)\cap s(C)\setminus r(f(A))\).
    Since \(d\in r(B\setminus f(A))\), we have
    \[
        r(f(A))\prec^{\bfd} d.
    \]
    On the other hand, since \(d\in s(C\setminus g(A))\), we have
    \[
        d\prec^{\bfd} s(g(A))=r(f(A)).
    \]
    This is impossible. Hence the amalgam is strong in the reduct, and \(\cK_0\) has the strong amalgamation property.

    Finally, we prove (iii) \(\Rightarrow\) (ii). Let \(\bfa,\bfb,\bfc\in\cK\), and let \(f\colon \bfa\hookrightarrow \bfb\) and \(g\colon \bfa\hookrightarrow \bfc\) be embeddings. Forgetting the order, we obtain embeddings \(f\colon \bfa_0\hookrightarrow \bfb_0\) and \(g\colon \bfa_0\hookrightarrow \bfc_0\). By strong amalgamation in \(\cK_0\), there exist \(\bfd_0\in\cK_0\) and embeddings \(r\colon \bfb_0\hookrightarrow \bfd_0\) and \(s\colon \bfc_0\hookrightarrow \bfd_0\)
    such that
    \[
        r\circ f=s\circ g
        \qquad\text{and}\qquad
        r(B)\cap s(C)=r(f(A)).
    \]
    Since the images of \(\bfb_0\) and \(\bfc_0\) intersect only in the common copy of \(\bfa_0\), and the orders on \(\bfb\) and \(\bfc\) agree there, the two induced linear orders on \(r(B)\subseteq D\) and \(s(C)\subseteq D\) can be amalgamated to a linear ordering \(\prec^{\bfd}\) of \(D\). If we now put \(\bfd=\langle \bfd_0,\prec^{\bfd}\rangle\), then \(\bfd\in\cK\) because \(\cK=\cK_0 * \cLO\), and the maps \(r\colon \bfb\hookrightarrow \bfd\) and \(s\colon \bfc\hookrightarrow \bfd\) are embeddings witnessing strong amalgamation in \(\cK\).
\end{proof}
The same argument gives the corresponding statement for strong joint embedding.

\begin{corollary}
    \label{cor:free-orders-sjep}
    Let \(L\), \(L_0\), \(\cK_0\), and \(\cK\) be as in Proposition~\ref{prop:free-orders-sap}. Then \(\cK\) has the strong joint embedding property if and only if \(\cK_0\) has the strong joint embedding property.
\end{corollary}

\begin{proof}
    Assume first that \(\cK\) has the strong joint embedding property. Let \(\bfb_0,\bfc_0\in\cK_0\). Choose arbitrary linear orderings \(\prec^{\bfb}\) and \(\prec^{\bfc}\) on \(B\) and \(C\), respectively, and form \(\bfb=\langle \bfb_0,\prec^{\bfb}\rangle\) and \(\bfc=\langle \bfc_0,\prec^{\bfc}\rangle\) in \(\cK\). By strong joint embedding in \(\cK\), there exist \(\bfd\in\cK\) and embeddings \(r\colon \bfb\hookrightarrow \bfd\) and \(s\colon \bfc\hookrightarrow \bfd\)
    with disjoint images. Passing to reducts shows that \(\cK_0\) has the strong joint embedding property.

    Conversely, assume that \(\cK_0\) has the strong joint embedding property. Let \(\bfb,\bfc\in\cK\). Then there exist \(\bfd_0\in\cK_0\) and embeddings \(r\colon \bfb_0\hookrightarrow \bfd_0\) and \(s\colon \bfc_0\hookrightarrow \bfd_0\)
    with disjoint images. Since the images are disjoint, the linear orders on \(\bfb\) and \(\bfc\) may be combined into a linear ordering \(\prec^{\bfd}\) of \(D\). Setting \(\bfd=\langle \bfd_0,\prec^{\bfd}\rangle\), we have \(\bfd\in\cK\), and the embeddings \(r\colon \bfb\hookrightarrow \bfd\) and \(s\colon \bfc\hookrightarrow \bfd\) have disjoint images. Thus \(\cK\) has the strong joint embedding property.
\end{proof}
In particular, if \(\cK_0\) is a Fraïssé class with the strong amalgamation property, then \(\cK_0 * \cLO\) is a reasonable Fraïssé order class. Many of the ordered graph, hypergraph, and metric-space classes considered earlier arise in exactly this way.

\section{Order-forgetful expansions and Ramsey transfer}

We conclude with a special situation in which the passage between an ordered class and its reduct is particularly simple. Roughly speaking, this happens when the additional order does not create new isomorphism types beyond those already present in the reduct.

\begin{definition}
    \label{def:order-forgetful}
    Let \(L\supseteq\{<\}\) be a signature, let \(L_0=L\setminus\{<\}\), and let \(\cK\) be a class of finite \(L\)-structures. We say that \(\cK\) is \emph{order-forgetful} if whenever \(\bfa,\bfb\in\cK\) satisfy \(\bfa_0\cong \bfb_0\), it follows that \(\bfa\cong \bfb\).
    Equivalently, for every isomorphism type in the reduct class \(\cK_0\), there is at most one isomorphism type in \(\cK\) lying above it.
\end{definition}

Thus the added order carries no additional isomorphism-type information.
Naturally ordered finite-dimensional vector spaces over a fixed finite field
are a basic example: once the underlying vector space is fixed, all natural
orderings are isomorphic.

For hereditary order-forgetful classes, the Ramsey property passes directly between the ordered and unordered settings.

\begin{proposition}
    \label{prop:order-forgetful-ramsey}
    Let \(L\supseteq\{<\}\) be a signature, let \(L_0=L\setminus\{<\}\), and let \(\cK\) be a hereditary order-forgetful class of finite \(L\)-structures. Put \(\cK_0=\{\bfa_0 : \bfa\in\cK\}\).
    Then \(\cK\) has the Ramsey property if and only if \(\cK_0\) has the Ramsey property.
\end{proposition}

\begin{proof}
    Assume first that \(\cK\) has the Ramsey property. Let \(\bfa_0\hookrightarrow \bfb_0\) be structures in \(\cK_0\), and let \(k\geq 2\). Choose an expansion \(\bfb\in\cK\) with reduct \(\bfb_0\). Since \(\bfa_0\hookrightarrow \bfb_0\), there is a copy \(\bfa_0'\leq \bfb_0\) of \(\bfa_0\). Let \(\bfa'\) be the ordered structure induced on \(\bfa_0'\) by the order of \(\bfb\). Because \(\cK\) is hereditary, we have \(\bfa'\in\cK\).
    Since \(\cK\) is order-forgetful and \(\bfa_0'\cong \bfa_0\), every expansion of \(\bfa_0\) in \(\cK\) is isomorphic to \(\bfa'\). In particular, if \(\bfa\in\cK\) is any expansion of \(\bfa_0\), then \(\bfa\cong \bfa'\leq \bfb\).
    Hence, replacing \(\bfa\) by an isomorphic copy if necessary, we may apply the
    Ramsey property of \(\cK\) to the pair \(\bfa\hookrightarrow \bfb\). Choose
    \(\bfc\in\cK\) such that
    \[
        \bfc\to (\bfb)^{\bfa}_k,
    \]
    and let \(\bfc_0=\bfc\!\restriction\!L_0\). We claim that
    \[
        \bfc_0\to (\bfb_0)^{\bfa_0}_k.
    \]

    Let
    \(c_0\colon \begin{pmatrix} \bfc_0 \\ \bfa_0 \end{pmatrix}\to\{1,\dots,k\}\)
    be a colouring. Every copy \(\bfa_0'\leq \bfc_0\) inherits from \(\bfc\) a linear ordering, and since \(\cK\) is hereditary, the induced ordered structure on \(\bfa_0'\) belongs to \(\cK\). Because \(\cK\) is order-forgetful and \(\bfa_0'\cong \bfa_0\), this induced ordered structure is isomorphic to \(\bfa\). Hence \(c_0\) induces a colouring
    \(c\colon \begin{pmatrix} \bfc \\ \bfa \end{pmatrix}\to\{1,\dots,k\}\)
    by setting \(c(\bfa')=c_0(\bfa_0')\),
    where \(\bfa_0'\) is the underlying reduct of the copy \(\bfa'\).

    Choose \(\bfb'\in \begin{pmatrix} \bfc \\ \bfb \end{pmatrix}\)
    homogeneous for \(c\). Then the reduct \(\bfb_0'\) is a copy of \(\bfb_0\) in \(\bfc_0\), and every copy of \(\bfa_0\) inside \(\bfb_0'\) inherits an ordered structure isomorphic to \(\bfa\). Since \(\bfb'\) is homogeneous for \(c\), it follows that \(\bfb_0'\) is homogeneous for \(c_0\). Thus
    \[
        \bfc_0\to (\bfb_0)^{\bfa_0}_k,
    \]
    and \(\cK_0\) has the Ramsey property.

    Conversely, assume that \(\cK_0\) has the Ramsey property. Let \(\bfa\hookrightarrow \bfb\) be structures in \(\cK\), and let \(k\geq 2\). Since \(\cK_0\) has the Ramsey property, there exists \(\bfc_0\in\cK_0\)
    such that
    \[
        \bfc_0\to (\bfb_0)^{\bfa_0}_k.
    \]
    Choose any expansion \(\bfc\in\cK\) with reduct \(\bfc_0\).

    We claim that
    \[
        \bfc\to (\bfb)^{\bfa}_k.
    \]
    Let \(c\colon \begin{pmatrix} \bfc \\ \bfa \end{pmatrix}\to\{1,\dots,k\}\)
    be a colouring. Define
    \(c_0\colon \begin{pmatrix} \bfc_0 \\ \bfa_0 \end{pmatrix}\to\{1,\dots,k\}\)
    by \(c_0(\bfa_0')=c(\bfa')\),
    where \(\bfa'\) is the ordered structure induced on \(\bfa_0'\) by the order of \(\bfc\). This is well defined because \(\cK\) is hereditary, and \(\bfa'\) is isomorphic to \(\bfa\) by order-forgetfulness.

    Choose \(\bfb_0'\in \begin{pmatrix} \bfc_0 \\ \bfb_0 \end{pmatrix}\)
    homogeneous for \(c_0\). Let \(\bfb'\) be the ordered structure induced on \(\bfb_0'\) by the order of \(\bfc\). Again, since \(\cK\) is hereditary and order-forgetful, we have \(\bfb'\cong \bfb\).
    Moreover, every copy of \(\bfa\) inside \(\bfb'\) has the same \(c\)-colour, because the corresponding copies of \(\bfa_0\) inside \(\bfb_0'\) all have the same \(c_0\)-colour. Hence \(\bfb'\) is homogeneous for \(c\), and therefore
    \[
        \bfc\to (\bfb)^{\bfa}_k.
    \]
    Thus \(\cK\) has the Ramsey property.
\end{proof}
The results of this chapter provide the structural link between an ordered Fraïssé class and its reduct. In the next chapter we use this link to study spaces of admissible orderings and to formulate the additional conditions under which ordered expansions determine the universal minimal flow.
